
\section{Wartości główne, jądra całek osobliwych i mnożniki Fouriera}
\label{sec:dt-singular}
\providecommand{\atlasref}[2]{#2}
Por. \atlasref{PDE-07::thm:pde7-laplace}{RRC-07: rozwiązanie fundamentalne laplasjanu}
oraz \atlasref{PDE-16::lem:pde16-distr-kernel}{RRC-16: jądra osobliwe}.

\subsection{Mnożniki Fouriera w przestrzeni $L^2$}
\label{sec:dt-multipliers-l2}

W tym rozdziale zachowujemy unitarną normalizację
transformacji Fouriera z definicji \ref{def:fourier-schwartz17}.
Materiał przenoszony z RRC-08 wymaga w tej konwencji
usunięcia czynników $(2\pi)^{-d/2}$ i $(2\pi)^{-d}$ w
tożsamościach norm oraz iloczynów skalarnych.
Por. \atlasref{PDE-08::thm:pde8-multiplier}{RRC-08: mnożniki Fouriera}.


\begin{twierdzenie}[Mnożnik w $L^2$]
\label{thm:dt-multiplier-l2}
Niech $m:\R^d\to\C$ będzie funkcją mierzalną i
istotnie ograniczoną. Dla $f\in L^2(\R^d)$ istnieje
dokładnie jedna klasa $T_mf\in L^2(\R^d)$ spełniająca
$\widehat{T_mf}=m\widehat f$ w sensie $L^2$.
Operator $T_m$ jest liniowy i ograniczony oraz
\[
 \|T_m\|_{L^2\to L^2}=\|m\|_{L^\infty}.
\]
Odwrotnie, jeśli $m$ jest mierzalna i skończona
prawie wszędzie, a ograniczony operator liniowy
$A:L^2\to L^2$ spełnia $\widehat{Af}=m\widehat f$
dla każdego $f\in L^2$ takiego, że
$m\widehat f\in L^2$, to $m\in L^\infty$
oraz $A=T_m$.
\end{twierdzenie}
\begin{proof}
Niech $M=\|m\|_{L^\infty}$. Z definicji istotnego
supremum dla każdego $k\ge1$ zbiór
$\{|m|>M+1/k\}$ ma miarę zero.
Ich przeliczalna suma jest zbiorem $\{|m|>M\}$,
więc $|m|\le M$ prawie wszędzie.
Dla $h\in L^2$ funkcja $mh$ jest więc kwadratowalna
i spełnia $\|mh\|_2\le M\|h\|_2$.
Mnożenie nie zależy od wyboru reprezentantów
klas funkcji $m$ i $h$.
Z twierdzenia~\ref{thm:plancherel22}
transformata Fouriera na $L^2$ jest bijekcją.
Możemy zatem określić
\[
 T_mf=\mathcal F^{-1}(m\widehat f).
\]
Jest to jedyna klasa spełniająca warunek
z tezy, a liniowość wynika z liniowości
wszystkich trzech operatorów w tym złożeniu.
Stosując dwukrotnie twierdzenie Plancherela,
otrzymujemy
\[
 \|T_mf\|_2=\|m\widehat f\|_2
 \le M\|\widehat f\|_2
 =M\|f\|_2.
\]
Stąd $\|T_m\|\le M$.

Gdy $M=0$, funkcja $m$ jest zerowa prawie
wszędzie i $T_m=0$, co daje równość norm.
Załóżmy $M>0$ i ustalmy $0<\varepsilon<M$.
Z definicji istotnego supremum zbiór
$E_\varepsilon=\{\xi:|m(\xi)|>M-\varepsilon\}$
ma dodatnią miarę: gdyby miał miarę zero,
$M-\varepsilon$ byłoby istotnym ograniczeniem
z góry dla $|m|$, wbrew definicji $M$.
Ponieważ $\R^d=\bigcup_{n\ge1}B_n(0)$,
dla pewnego $n$ zbiór
$E=E_\varepsilon\cap B_n(0)$ ma dodatnią
i skończoną miarę. Niech
\[
 h=\frac{\chi_E}{\lambda_d(E)^{1/2}},
 \qquad f=\mathcal F^{-1}h.
\]
Wówczas $\|h\|_2=1$ i
$\|f\|_2=1$. Ponadto
\[
 \|mh\|_2^2
 =\frac1{\lambda_d(E)}\int_E|m(\xi)|^2
 \,\mathrm d\lambda_d(\xi)
 \ge(M-\varepsilon)^2.
\]
W konsekwencji
$\|T_mf\|_2/\|f\|_2=\|mh\|_2\ge M-\varepsilon$.
Przechodząc z $\varepsilon$ do zera,
otrzymujemy $\|T_m\|\ge M$.

Dla implikacji odwrotnej połóżmy $C=\|A\|$.
Jeśli zbiór $\{|m|>C+\varepsilon\}$ dla pewnego
$\varepsilon>0$ ma dodatnią miarę, to z
prawie wszędzie skończoności $m$ wynika, że
z dokładnością do zbioru miary zero jest on
przeliczalną sumą poniższych zbiorów po $n\ge1$.
Dodatnia miara sumy oznacza, że dla pewnego $n$ zbiór
\[
 E=\{\xi:C+\varepsilon<|m(\xi)|\le n\}
       \cap B_n(0)
\]
ma dodatnią miarę. Ma też miarę skończoną.
Wybierając $f=\mathcal F^{-1}\chi_E$,
mamy $m\widehat f\in L^2$ dzięki ograniczeniu
$|m|\le n$ na $E$. Z założenia i Plancherela
wynika wtedy
\[
 \frac{\|Af\|_2}{\|f\|_2}
 =\frac{\|m\chi_E\|_2}{\|\chi_E\|_2}
 \ge C+\varepsilon>C,
\]
co przeczy definicji normy operatorowej.
Zatem dla każdego $\varepsilon>0$ zbiór
$\{|m|>C+\varepsilon\}$ ma miarę zero.
Biorąc przeliczalną sumę dla
$\varepsilon=1/k$, $k\ge1$, dostajemy
$|m|\le C$ prawie wszędzie. Tym samym
$m\in L^\infty$, a warunek na $A$ obowiązuje
już dla wszystkich $f\in L^2$.
Jednoznaczność daje $A=T_m$.
\end{proof}

\begin{wniosek}[Rachunek ograniczonych mnożników]
Dla $m,n\in L^\infty(\R^d)$ i $a,b\in\C$ zachodzi
\[
 T_{am+bn}=aT_m+bT_n,\qquad
 T_mT_n=T_{mn},\qquad T_1=I,\qquad
 T_m^*=T_{\overline m}.
\]
Gwiazdka oznacza operator sprzężony w zespolonej
przestrzeni Hilberta $L^2$.
\end{wniosek}
\begin{proof}
Dla $f\in L^2$ z definicji mnożników mamy
\[
 \begin{aligned}
 \mathcal F(aT_mf+bT_nf)&=(am+bn)\widehat f,\\
 \mathcal F(T_mT_nf)&=m\mathcal F(T_nf)=mn\widehat f,\\
 \mathcal F(T_1f)&=\widehat f.
 \end{aligned}
\]
Wszystkie symbole po prawej stronie są istotnie
ograniczone. Równości operatorów wynikają
z iniektywności transformaty na $L^2$.
Dla $f,g\in L^2$ tożsamość
 z twierdzenia~\ref{thm:plancherel22} daje
\[
 \begin{aligned}
 \langle T_m f,g\rangle_2
 &=\int_{\R^d}m(\xi)\widehat f(\xi)
       \overline{\widehat g(\xi)}\,\mathrm d\lambda_d(\xi)\\
 &=\int_{\R^d}\widehat f(\xi)
       \overline{\overline{m(\xi)}\widehat g(\xi)}
       \,\mathrm d\lambda_d(\xi)\\
 &=\langle f,T_{\overline m}g\rangle_2.
 \end{aligned}
\]
Całki są bezwzględnie zbieżne przez nierówność
Cauchy'ego--Schwarza i istotną ograniczoność $m$.
To dowodzi równości operatorów sprzężonych.
\end{proof}

\begin{przyklad}[Transformata Riesza w $L^2$]
Dla $1\le j\le d$ połóżmy
$m_j(\xi)=-i\xi_j/|\xi|$ dla $\xi\ne0$ i $m_j(0)=0$.
Twierdzenie~\ref{thm:dt-multiplier-l2} definiuje
ograniczony operator $R_j=T_{m_j}:L^2\to L^2$,
zwany transformatą Riesza.
Mamy $|m_j(\xi)|\le1$. Dla każdego
$0<\varepsilon<1$ zbiór
$\{\xi\ne0:|m_j(\xi)|>1-\varepsilon\}$
zawiera otwartą kulę o dodatnim promieniu
wokół $e_j$: funkcja $|\xi_j|/|\xi|$ jest
ciągła w otoczeniu $e_j$ i przyjmuje tam
wartość $1$ w punkcie $e_j$.
Zatem $\|m_j\|_\infty=1$ i w konsekwencji
$\|R_j\|_{L^2\to L^2}=1$, także dla $d=1$.

Ponieważ $\overline{m_j}=-m_j$ oraz dla
$\xi\ne0$ zachodzą
\[
 \sum_{j=1}^d |m_j(\xi)|^2=1,\qquad
 \sum_{j=1}^d m_j(\xi)^2=-1,
\]
z rachunku mnożników otrzymujemy
\[
 R_j^*=-R_j,\qquad
 \sum_{j=1}^d R_j^2=-I,\qquad
 \sum_{j=1}^d\|R_j f\|_2^2=\|f\|_2^2.
\]
Ostatnia tożsamość wynika również bezpośrednio
z Plancherela:
\[
 \sum_{j=1}^d\|R_j f\|_2^2
 =\int_{\R^d}
 \left(\sum_{j=1}^d|m_j(\xi)|^2\right)
 |\widehat f(\xi)|^2\,\mathrm d\lambda_d(\xi)
 =\|f\|_2^2.
\]
Wartość mnożnika w pojedynczym punkcie $\xi=0$
nie zmienia operatora na $L^2$, ponieważ
$\lambda_d(\{0\})=0$. Nie oznacza to jednak,
że można przez ten sam symbol mnożyć dowolną
dystrybucję temperowaną: wzdłuż $\xi=te_j$
symbol przyjmuje wartość $-i$ dla $t>0$
oraz $i$ dla $t<0$. Nie ma więc ciągłego
przedłużenia do zera i nie spełnia warunku
mnożenia dystrybucji z poprzedniej sekcji.
Na przykład dla testu $\varphi$ z $\varphi(0)\ne0$
iloczyn $m_j\varphi$ nie jest ciągły w zerze,
a zatem nie należy do $\mathcal S$.
\end{przyklad}

\begin{uwaga}[Zakres twierdzenia o mnożnikach]
Kryterium istotnej ograniczoności symbolu
dotyczy normy $L^2$ i korzysta z twierdzenia
Plancherela. Dla $p\ne2$ powyższy dowód
nie daje twierdzenia o ograniczoności
na $L^p$. W szczególności nie należy
wyciągać z niego wniosku o ograniczoności
transformat Riesza w normach $L^1$ lub
$L^\infty$. Badanie mnożników na $L^p$
dla $1<p<\infty$ wymaga dalszych narzędzi
analizy harmonicznej; do tej teorii wrócimy po omówieniu operatorów maksymalnych i interpolacji.
\end{uwaga}

\subsection{Jądra drugich pochodnych potencjału Newtona}
Niech $\omega_{d-1}=\mathcal H^{d-1}(S^{d-1})$ oznacza pole
jednostkowej sfery, a $\Gamma$ rozwiązanie fundamentalne $-\Delta$:
$-\Delta\Gamma=\delta_0$. Z twierdzenia~\ref{thm:dt-laplace} dla $d\ge3$ można wziąć
$\Gamma(x)=\bigl((d-2)\omega_{d-1}\bigr)^{-1}|x|^{2-d}$,
a dla $d=2$ funkcję $\Gamma(x)=-(2\pi)^{-1}\log|x|$.
Dla $x\ne0$ w każdym z obu przypadków
$\partial_i\Gamma(x)=-x_i/(\omega_{d-1}|x|^d)$.
Różniczkując względem $x_j$, otrzymujemy
\begin{equation}\label{eq:dt-cz-kernel}
 K_{ij}(x):=\partial_{ij}\Gamma(x)
 =\frac{d x_ix_j-\delta_{ij}|x|^2}
 {\omega_{d-1}|x|^{d+2}},\qquad 1\le i,j\le d,
\end{equation}
gdzie $\delta_{ij}$ oznacza deltę Kroneckera.
Jądro jest gładkie poza zerem i spełnia
$K_{ij}(tx)=t^{-d}K_{ij}(x)$ dla $t>0$.
Różniczkowanie tej tożsamości daje jednorodność gradientu
stopnia $-d-1$. Z ograniczoności obu funkcji na zwartej
sferze wynika $|K_{ij}(x)|\le C_d|x|^{-d}$ oraz
$|\nabla K_{ij}(x)|\le C_d|x|^{-d-1}$.
Dla $i\ne j$ odbicie względem $\theta_i=0$ zmienia znak
$\theta_i\theta_j$ i zachowuje miarę powierzchniową.
Dla $i=j$ całki z $\theta_i^2$ są równe przez symetrię
permutacji współrzędnych, a ich suma wynosi $\omega_{d-1}$.
Stąd
\[
 \int_{S^{d-1}}\theta_i\theta_j\,\mathrm d\mathcal H^{d-1}
 =\frac{\omega_{d-1}}d\delta_{ij},\qquad
 \int_{S^{d-1}}K_{ij}(\theta)\,\mathrm d\mathcal H^{d-1}(\theta)=0.
\]
Każde jądro jest niezerowe na pewnym otwartym podzbiorze
sfery (dla $i=j$ weźmy $\theta=e_i$, dla $i\ne j$
$\theta=(e_i+e_j)/\sqrt2$). Współrzędne biegunowe dają
\[
 \int_{B_1(0)}|K_{ij}(z)|\,\mathrm d\lambda_d(z)
 =\int_{S^{d-1}}|K_{ij}(\theta)|\,\mathrm d\mathcal H^{d-1}(\theta)
 \int_0^1 r^{-1}\,\mathrm d\lambda_1(r)=+\infty .
\]
Jeżeli $f$ jest ciągła i $f(x)\ne0$, to $|f(x-z)|$
jest oddzielone od zera w otoczeniu zera. Zwykła całka
Lebesgue'a z $K_{ij}(z)f(x-z)$ nie jest więc bezwzględnie
zbieżna, nawet dla takich $f\in C_c^\infty$.

\subsection{Główna wartość całki i dystrybucja związana z jądrem}

\begin{definicja}[Główna wartość całki]
\label{def:dt-principal-value}
Niech $f\in C_c^1(\R^d)$, a $K$ będzie jądrem klasy $C^1$ poza
zerem, jednorodnym stopnia $-d$ i o zerowej średniej sferycznej.
Definiujemy
\[
 T_Kf(x)=\operatorname{p.v.}\!\int_{\R^d}K(x-y)f(y)\,\mathrm d\lambda_d(y)
 :=\lim_{\varepsilon\downarrow0}
 \int_{|x-y|>\varepsilon}K(x-y)f(y)\,\mathrm d\lambda_d(y).
\]
Dla $0<\varepsilon<1$ podstawiamy $z=x-y$.
Jednorodność i zerowa średnia dają zerową całkę z $K$
po każdym pierścieniu. Zatem po odjęciu $f(x)$
w pierścieniu $\varepsilon<|z|<1$ dostajemy w granicy
\[
 \begin{aligned}
 T_Kf(x)&=\int_{|z|<1}K(z)(f(x-z)-f(x))\,\mathrm d\lambda_d(z)\\
 &\quad+\int_{|z|\ge1}K(z)f(x-z)\,\mathrm d\lambda_d(z).
 \end{aligned}
\]
Obie całki są bezwzględnie zbieżne: w pierwszej
$|f(x-z)-f(x)|\le\|\nabla f\|_\infty |z|$,
więc całka radialna jest ograniczona przez stałą razy
$\int_0^1\mathrm d\lambda_1(r)$; druga ma moduł
nie większy niż $C\|f\|_1$. Uzasadnia to istnienie
granicy dla każdego $x$ oraz
$|T_Kf(x)|\le C(\|\nabla f\|_\infty+\|f\|_1)$.
\end{definicja}

Dystrybucję $\operatorname{p.v.}K$ określamy przez
\[
 \begin{aligned}
 \langle\operatorname{p.v.}K,\varphi\rangle
 &=\int_{|z|<1}K(z)(\varphi(z)-\varphi(0))\,\mathrm d\lambda_d(z)\\
 &\quad+\int_{|z|\ge1}K(z)\varphi(z)\,\mathrm d\lambda_d(z).
 \end{aligned}
\]
Dla testów o nośniku w ustalonym zwartym $L$ moduł
działania jest ograniczony przez
$C\|\nabla\varphi\|_\infty+C\lambda_d(L)\|\varphi\|_\infty$.
Liniowość i oszacowanie przez półnormy na każdym zwartym
nośniku dowodzą ciągłości dystrybucyjnej na mocy
twierdzenia \ref{thm:distribution-estimate19}.
Jest to także dystrybucja temperowana. Dla
$\varphi\in\mathcal S$ pierwsza całka jest ograniczona
przez stałą razy $\|\nabla\varphi\|_\infty$,
a druga przez
$C Q_{d+1,0}(\varphi)\int_1^\infty r^{-1}(1+r)^{-d-1}\,\mathrm d\lambda_1(r)$,
z dodatkowym stałym czynnikiem powierzchni sfery.
To daje ciągłość na $\mathcal S$.

\begin{lemat}[Hessian rozwiązania fundamentalnego jako dystrybucja]
\label{lem:dt-cz-distr-kernel}
Dla każdego $1\le i,j\le d$ zachodzi w $\mathcal D'(\R^d)$
\begin{equation}\label{eq:dt-cz-distr-kernel}
 \partial_{ij}\Gamma
 =\operatorname{p.v.}K_{ij}-\frac{\delta_{ij}}d\delta_0.
\end{equation}
W szczególności dla $f\in C_c^\infty(\R^d)$ i $u=\Gamma*f$,
\begin{equation}\label{eq:dt-cz-repr}
 \partial_{ij}u=T_{ij}f-\frac{\delta_{ij}}d f,
 \qquad T_{ij}=T_{K_{ij}}.
\end{equation}
\end{lemat}
\begin{proof}
Niech $\varphi\in C_c^\infty(\R^d)$, $R>1$ tak duże,
aby nośnik $\varphi$ leżał w $B_R(0)$, oraz
$D_\varepsilon=\{x:\varepsilon<|x|<R\}$.
Z lokalnej całkowalności $\Gamma$ działanie
$\partial_{ij}\Gamma$ jest granicą całki
$\int_{D_\varepsilon}\Gamma\partial_{ij}\varphi\,\mathrm d\lambda_d$.
Całkując przez części najpierw względem $x_i$, potem
$x_j$, otrzymujemy
\[
 \begin{aligned}
 \int_{D_\varepsilon}\Gamma\partial_{ij}\varphi\,\mathrm d\lambda_d
 &=\int_{D_\varepsilon}K_{ij}\varphi\,\mathrm d\lambda_d\\
 &\quad+\int_{\partial D_\varepsilon}
 \Gamma n_i\partial_j\varphi\,\mathrm d\mathcal H^{d-1}\\
 &\quad-\int_{\partial D_\varepsilon}
 (\partial_i\Gamma)n_j\varphi\,\mathrm d\mathcal H^{d-1}.
 \end{aligned}
\]
Na zewnętrznej sferze oba wyrazy brzegowe znikają.
Na wewnętrznej sferze $x=\varepsilon\theta$ mamy
$n=-\theta$, $\partial_i\Gamma(x)=
-\theta_i/(\omega_{d-1}\varepsilon^{d-1})$,
a miara powierzchniowa skaluje się przez
$\varepsilon^{d-1}$. Ostatni wyraz wynosi więc
\[
 -\frac1{\omega_{d-1}}\int_{S^{d-1}}\theta_i\theta_j
 \varphi(\varepsilon\theta)\,\mathrm d\mathcal H^{d-1}(\theta)
 \longrightarrow-\frac{\delta_{ij}}d\varphi(0).
\]
Przedostatni wyraz jest ograniczony co do modułu przez
$C\|\partial_j\varphi\|_\infty\varepsilon|\log\varepsilon|$
dla $d=2$ i przez $C_d\|\partial_j\varphi\|_\infty\varepsilon$
dla $d\ge3$, więc znika w granicy.
Pierwsza całka dąży do działania $\operatorname{p.v.}K_{ij}$.
Daje to~\eqref{eq:dt-cz-distr-kernel}.
Z twierdzeń~\ref{thm:smoothing21} i
\ref{thm:convolution-derivative21} splot
$u=\Gamma*f$ jest gładki, $-\Delta u=\delta_0*f=f$ oraz
$\partial_{ij}u=(\partial_{ij}\Gamma)*f$.
Działanie $\operatorname{p.v.}K_{ij}$ na teście
$z\mapsto f(x-z)$ jest dokładnie $T_{ij}f(x)$,
a $\delta_0*f=f$. Otrzymujemy~\eqref{eq:dt-cz-repr}.
\end{proof}

\subsection{Mnożniki związane z jądrem Newtona}
\begin{lemat}[Oszacowanie w $L^2$]
\label{lem:dt-cz-l2}
Dla $f\in C_c^\infty(\R^d)$ zachodzi
$\|T_{ij}f\|_2\le C_d\|f\|_2$. Operator $T_{ij}$
przedłuża się jednoznacznie do ograniczonego operatora na
$L^2(\R^d)$.
\end{lemat}
\begin{proof}
Najpierw wykażemy przynależność do $L^2$, bez jednostajnego
oszacowania normy przez $\|f\|_2$. Z definicji głównej
wartości $T_{ij}f$ jest ograniczona na całej przestrzeni.
Niech $\operatorname{supp}f\subset B_R(0)$, $R\ge1$.
Dla $|x|>2R$ mamy $|x-y|\ge|x|/2$ na nośniku $f$,
więc całka nie jest osobliwa i
$|T_{ij}f(x)|\le C_d|x|^{-d}\|f\|_1$.
Całka kwadratu tego ograniczenia poza $B_{2R}$ jest
skończona, gdyż $\int_{2R}^\infty r^{-d-1}
\,\mathrm d\lambda_1(r)<\infty$.
Na $B_{2R}$ wystarcza ograniczoność. Stąd $T_{ij}f\in L^2$,
a z~\eqref{eq:dt-cz-repr} także $H_{ij}:=\partial_{ij}u\in L^2$.

Z $-\Delta u=f$ i przemienności pochodnych mamy
$-\Delta H_{ij}=\partial_{ij}f$ dystrybucyjnie.
Obie strony są dystrybucjami temperowanymi; równość
w $\mathcal S'$ wynika z gęstości $C_c^\infty$ w $\mathcal S$
(twierdzenie~\ref{thm:density-schwartz17}). Twierdzenie~\ref{thm:fourier-auto-tempered22} daje
\[
 |\xi|^2\widehat H_{ij}=-\xi_i\xi_j\widehat f
 \qquad\text{w }\mathcal S'(\R^d).
\]
Obie strony są funkcjami lokalnie całkowalnymi, ponieważ
$\widehat H_{ij},\widehat f\in L^2$ z twierdzenia
Plancherela. Z jednoznaczności dystrybucji regularnych
(stwierdzenie~\ref{prop:regular-injective19}) mamy równość
prawie wszędzie. Dzieląc dla $\xi\ne0$, otrzymujemy
$\widehat H_{ij}(\xi)=-(\xi_i\xi_j/|\xi|^2)\widehat f(\xi)$
prawie wszędzie. Nie ma dodatkowego składnika skupionego
w zerze, bo $\widehat H_{ij}$ jest funkcją $L^2$.
Z~\eqref{eq:dt-cz-repr} wynika
$\widehat{T_{ij}f}=m_{ij}\widehat f$, gdzie
$m_{ij}(\xi)=\delta_{ij}/d-\xi_i\xi_j/|\xi|^2$
poza zerem, a wartość w zerze jest nieistotna.
Mamy $|m_{ij}|\le1+1/d\le2$. Z twierdzenia
~\ref{thm:plancherel22} przy przyjętej konwencji
\[
 \|T_{ij}f\|_2=\|m_{ij}\widehat f\|_2
 \le 2\|\widehat f\|_2=2\|f\|_2.
\]
Argument działa również dla $d=2$, bez wymagania, aby
sam potencjał $u$ należał do $L^2$. Liniowość, gęstość
$C_c^\infty$ w $L^2$ i zupełność $L^2$ dają jednoznaczne
ciągłe przedłużenie (gęstość w $L^2$ jest wykorzystana w dowodzie
twierdzenia~\ref{thm:plancherel22}).
\end{proof}
